\section{配对、准粒子与阻塞}\label{sec:2-3}


在低能核物理中，原子核不能简单视作各核子在静态平均场中独立运动的费米气体体系。由于核子—核子有效相互作用中短程吸引成分的存在，核子体系展现出显著的两体关联。在实验经验上，近稳定线已测量的绝大多数偶偶核基态均表现为自旋与宇称 $I^\pi = 0^+$，且其内禀破对激发能级通常明显高于邻近奇 $A$ 核的低能单粒子激发态；与此同时，结合能系统学呈现出普遍的奇偶交错效应。这些现象表明，同类核子之间存在强烈的超流配对倾向。为描述这种超越独立粒子平均场的集体效应，通常引入基于 Bardeen–Cooper–Schrieffer（BCS）理论的配对均场与准粒子方法\cite{book-suhonen-from-nucleons-to-nucleus, book-jain-nuclear-isomers}。

将配对微观理论应用于有限原子核体系，需要处理时间反演对称性、离散单粒子能级结构以及粒子数规范对称性破缺等一系列理论问题。本节系统阐述配对、准粒子与阻塞效应的微观机制：首先从时间反演轨道与短程吸引相互作用出发，分析偶偶核基态配对与奇偶质量差的微观成因，明确经验质量差与纯配对相互作用参数的物理边界；随后在平均粒子数约束变分框架下，给出 Bogoliubov–Valatin 正则变换、准粒子哈密顿量的 $2\times 2$ 矩阵化与对角化，讨论占据振幅、准粒子能量、自洽能隙方程与粒子数方程；进而辨析配对能隙、配对张量振幅与规范角三者的物理含义，说明均场粒子数破缺的理论属性，引出局域配对转动与粒子数投影接口；随后分析奇核子引起的泡利阻塞与自洽均场重排，探讨多准粒子态的形成机制并澄清静态 BCS 解在弱耦合下均场行为的物理本质；最后结合简化的两能级配对模型分析其临界解析行为与适用条件，阐明微观配对对集体动力学响应的调制作用。



\subsection{时间反演成对、短程吸引与奇偶效应的微观根源}\label{sec:2-3-1}


时间反演不变的平均场使半整数核子的单粒子轨道具有Kramers伙伴。先以球形场中的 $|k\rangle=|a,j,m_k\rangle$ 为例，a包含其它量子数；在轴对称变形场中，相应伙伴仍可定义，但j通常不再是好量子数。在常用的 Condon–Shortley 相位约定下\cite{book-suhonen-from-nucleons-to-nucleus}：
\begin{equation}\label{eq:2-3-1}
|\bar{k}\rangle \equiv \mathcal{T} |k\rangle = (-1)^{j - m_k} |a, -m_k\rangle.
\end{equation}
由于核子是自旋 $1/2$ 的费米子，重复作用时间反演算符满足 $\mathcal{T}^2 = -1$，保证了单粒子能级在无外加时间反演破缺场时至少具有二重 Kramers 简并。

当两个处于相互时间反演轨道的同类核子（质子—质子或中子—中子）发生相互作用时，在特定磁子能级对上产生对态 $c_k^\dagger c_{\bar{k}}^\dagger |0\rangle$。需要明确区分单个磁轨道对态与全壳层角动量耦合态的物理表示：单个固定磁投影的轨道对态 $c_k^\dagger c_{\bar{k}}^\dagger |0\rangle$ 仅具有总磁投影 $M = m_k + (-m_k) = 0$，在空间三维旋转下并非单一总角动量 $J$ 的本征态，而是包含了允许的总角动量 $J = 0, 2, 4, \dots$ 的线性叠加。在球对称单 $j$ 壳层中，将同壳层核子对严格耦合为总角动量为零的球形标量态时，需要对该壳层内所有磁子能级进行角动量张量耦合求和。在真空归一化约定下，归一化的双核子 $J=0$ 态为：
\begin{equation}\label{eq:2-3-2}
\begin{aligned}
|j^2; J=0, M=0\rangle &= \frac{1}{\sqrt{2(2j+1)}} \sum_{m=-j}^j (-1)^{j-m} c_{jm}^\dagger c_{j,-m}^\dagger |0\rangle \\
&= \frac{1}{\sqrt{\Omega}} \sum_{m>0} (-1)^{j-m} c_{jm}^\dagger c_{j,-m}^\dagger |0\rangle,
\end{aligned}
\end{equation}
其中 $\Omega = (2j+1)/2$ 为该单 $j$ 壳层的 Kramers 轨道对简并度。在轴对称形变核中，球形角动量 $J$ 被形变场破缺，但时间反演对称性依然保持，$c_k^\dagger c_{\bar{k}}^\dagger$ 对应于沿内禀对称轴投影 $K=0$ 的时间反演对产生算符。

反向运动的经典轨道可帮助形成成对的直觉，但不能单靠这幅图像证明任意两体作用的角动量排序。下面采用同类核子的吸引配对模型，主要保留时间反演轨道间的对散射。球形情况下，跨磁轨道的相干求和可以形成上述J=0对；实际相互作用还含其它角动量通道。实验奇偶效应为配对机制提供支持，具体强度仍取决于壳结构与所选相互作用\cite{book-suhonen-from-nucleons-to-nucleus, book-jain-nuclear-isomers}。

这种配对吸引直接反映在结合能的奇偶效应中。在偶偶核中，价核子以成对方式占据各能级；在奇 $A$ 核中，至少存在一个未配对的奇核子，该核子无法以同样方式形成时间反演基态对，导致奇 $A$ 核的总结合能系统性地低于相邻偶偶核的平均内插值；在奇奇核中，同时存在未成对的价质子与价中子，结合能进一步受损。在核质量实验数据分析中，常采用正束缚能 $B(N, Z) > 0$（结合越紧密数值越大）的离散差分公式来提取奇偶质量差。三点差分与四点差分公式定义为：
\begin{equation}\label{eq:2-3-3}
\Delta_B^{(3)}(N, Z) \equiv \frac{(-1)^{N+1}}{2} \left[ B(N-1, Z) - 2B(N, Z) + B(N+1, Z) \right],
\end{equation}
\begin{equation}\label{eq:2-3-4}
\Delta_B^{(4)}(N, Z) \equiv \frac{(-1)^N}{4} \left[ B(N-2, Z) - 3B(N-1, Z) + 3B(N, Z) - B(N+1, Z) \right].
\end{equation}
四点差分公式在结构上严格对应于相邻两个三点指标的对称中心平均：
\begin{equation}\label{eq:2-3-5}
\Delta_B^{(4)}(N, Z) = \frac{\Delta_B^{(3)}(N-1, Z) + \Delta_B^{(3)}(N, Z)}{2}.
\end{equation}
若设定束缚能包含平滑背景与纯奇偶交错项 $B(N) = B_{\mathrm{sm}}(N) + p(-1)^N$（$p > 0$），则上述带有 $(-1)^{N+1}$ 前置因子的三点式与四点式对纯交错项均给出一致的正能标 $+2p$（若改用总能量 $E = -B < 0$，则相应三点式前置因子调整为 $(-1)^N$）。

必须在物理上明确：\textbf{由实验结合能差分公式提取的经验奇偶质量差 $\Delta_B^{(3)}$ 或 $\Delta_B^{(4)}$，不能直接等同于微观哈密顿量中的纯配对能隙参数 $\Delta$。} 实验差分不仅包含配对关联能，还混杂了费米面附近离散单粒子能级间距 $(\epsilon_{k+1} - \epsilon_k)$、核子数变化引起的自洽均场重排能、形变平衡点的偏移以及单极势漂移等净贡献。在幻数闭壳邻近，差分主要反映单粒子能隙跳变；即使在开壳核区，将经验差分直接代入微观模型作为恒定配对参数也会忽略轨道依赖的能隙结构与极化响应。



\subsection{约束变分、Bogoliubov 变换与 BCS 均场推导}\label{sec:2-3-2}


为了在多体框架下统一处理配对关联，BCS 理论引入了破缺粒子数守恒的试探波函数。设单粒子产生算符为 $c_k^\dagger$，时间反演共轭算符为 $c_{\bar{k}}^\dagger$。为避免在配对求和中出现重复计数，指标 $k > 0$ 遍历半数单粒子态空间（例如在轴对称或球形基底下选取投影 $m_k > 0$）。定义单轨道对产生算符为 $P_k^\dagger \equiv c_k^\dagger c_{\bar{k}}^\dagger$。


\subsubsection{BCS 拟设与粒子数约束}\label{sec:2-3-bcs拟设与粒子数约束}


BCS 基态试探波函数假定时间反演轨道对处于独立相干凝聚态\cite{book-suhonen-from-nucleons-to-nucleus}：
\begin{equation}\label{eq:2-3-6}
|\mathrm{BCS}\rangle = \prod_{k>0} \left( u_k + v_k P_k^\dagger \right) |0\rangle = \prod_{k>0} \left( u_k + v_k c_k^\dagger c_{\bar{k}}^\dagger \right) |0\rangle,
\end{equation}
其中 $|0\rangle$ 为物理真空，$u_k$ 与 $v_k$ 为实变分参量。态矢的模方展开为 $\langle\mathrm{BCS}|\mathrm{BCS}\rangle = \prod_{k>0} (u_k^2 + v_k^2)$。乘积等于 1 本身在逻辑上并不能唯一推出每个因子独立等于 1；在理论中，我们将逐对归一化条件：
\begin{equation}\label{eq:2-3-7}
u_k^2 + v_k^2 = 1 \quad (\forall k > 0)
\end{equation}
确立为独立对振幅的标准规范化约定。在该约定下，$v_k^2$ 与 $u_k^2$ 分别代表轨道对 $(k, \bar{k})$ 被一对核子占据与未被占据的几率。

由于波函数展开包含具有不同偶数粒子数的态分量，$|\mathrm{BCS}\rangle$ 破缺了 $U(1)$ 粒子数规范对称性。为使该态对应特定核素的平均粒子数，需通过拉格朗日乘子 $\lambda$（化学势）引入变分约束：
\begin{equation}\label{eq:2-3-8}
\langle\mathrm{BCS}| \hat{n} |\mathrm{BCS}\rangle = \bar{n}, \quad \hat{n} = \sum_{k>0} \left( c_k^\dagger c_k + c_{\bar{k}}^\dagger c_{\bar{k}} \right).
\end{equation}
计算给出平均核子数关系式：
\begin{equation}\label{eq:2-3-9}
\bar{n} = 2 \sum_{k>0} v_k^2.
\end{equation}
式中因子 $2$ 源于每个轨道对包含两个时间反演态。需要强调的是，平均数约束 $\langle\hat{n}\rangle = \bar{n}$ 仅在期望值层面上控制粒子数，并不等于微观状态具有确定粒子数。


\subsubsection{Bogoliubov–Valatin 正则变换}\label{sec:2-3-bogoliubovvalatin正则变换}


定义准粒子产生与消灭算符如下\cite{book-suhonen-from-nucleons-to-nucleus}：
\begin{equation}\label{eq:2-3-10}
\alpha_k^\dagger = u_k c_k^\dagger - v_k c_{\bar{k}}, \quad \alpha_{\bar{k}}^\dagger = u_k c_{\bar{k}}^\dagger + v_k c_k,
\end{equation}
其厄米共轭为：
\begin{equation}\label{eq:2-3-11}
\alpha_k = u_k c_k - v_k c_{\bar{k}}^\dagger, \quad \alpha_{\bar{k}} = u_k c_{\bar{k}} + v_k c_k^\dagger.
\end{equation}
由逐对归一条件 $u_k^2 + v_k^2 = 1$ 可直接验证准粒子算符满足费米子反对易关系：
\begin{equation}\label{eq:2-3-12}
\{\alpha_k, \alpha_{k'}^\dagger\} = \delta_{kk'}, \quad \{\alpha_{\bar{k}}, \alpha_{\bar{k}'}^\dagger\} = \delta_{kk'}, \quad \{\alpha_k, \alpha_{k'}\} = 0.
\end{equation}
作用消灭算符于 BCS 基态可知 $\alpha_k |\mathrm{BCS}\rangle = 0$（$\forall k > 0$），因此 $|\mathrm{BCS}\rangle$ 即为无准粒子的准粒子真空态。


\subsubsection{辅助哈密顿量与 $2\times 2$ 矩阵结构}\label{sec:2-3-辅助哈密顿量与2times2矩阵结构}


考虑包含单粒子能量与对散射相互作用的模型纯配对哈密顿量：
\begin{equation}\label{eq:2-3-13}
\hat{H} = \sum_{k>0} \epsilon_k \left( c_k^\dagger c_k + c_{\bar{k}}^\dagger c_{\bar{k}} \right) - \sum_{k, l > 0} G_{kl} c_k^\dagger c_{\bar{k}}^\dagger c_{\bar{l}} c_l = \sum_{k>0} \epsilon_k \left( c_k^\dagger c_k + c_{\bar{k}}^\dagger c_{\bar{k}} \right) - \sum_{k, l > 0} G_{kl} P_k^\dagger P_l,
\end{equation}
其中相互作用矩阵元假定满足实对称 $G_{kl} = G_{lk} > 0$。引入平均粒子数约束后的辅助哈密顿量为 $\hat{H}' = \hat{H} - \lambda \hat{n}$。通过准粒子逆变换表达粒子算符：
\begin{equation}\label{eq:2-3-14}
c_k^\dagger = u_k \alpha_k^\dagger + v_k \alpha_{\bar{k}}, \quad c_{\bar{k}}^\dagger = u_k \alpha_{\bar{k}}^\dagger - v_k \alpha_k,
\end{equation}
并按准粒子正规积展开，辅助哈密顿量精确分解为：
\begin{equation}\label{eq:2-3-15}
\hat{H}' = H_0' + H_{11}' + H_{20}' + V_{\mathrm{RES}}.
\end{equation}
其中 $H_0' = \langle\mathrm{BCS}| \hat{H}' |\mathrm{BCS}\rangle$ 为基态平均能量；$H_{11}'$ 仅包含准粒子数算符项；$H_{20}'$ 包含成对产生或消灭准粒子的项；$V_{\mathrm{RES}}$ 包含四准粒子残余相互作用。

对哈密顿量中的配对两体相互作用项进行直接 Wick 收缩，保留完整的 $k=l$ 对角配对项及其平均场贡献，得到自洽自能修正：
\begin{equation}\label{eq:2-3-16}
\mu_k = G_{kk} v_k^2.
\end{equation}
进而定义相对于化学势与自能的单粒子有效能量：
\begin{equation}\label{eq:2-3-17}
\xi_k \equiv \epsilon_k - \lambda - \mu_k = \epsilon_k - \lambda - G_{kk} v_k^2,
\end{equation}
以及自洽配对能隙：
\begin{equation}\label{eq:2-3-18}
\Delta_k \equiv \sum_{l > 0} G_{kl} u_l v_l.
\end{equation}
在上述模型定义下，非对角的准粒子对产生项系数整齐表示为：
\begin{equation}\label{eq:2-3-19}
H_{20}' = \sum_{k>0} \left[ 2\xi_k u_k v_k - \Delta_k (u_k^2 - v_k^2) \right] \left( \alpha_k^\dagger \alpha_{\bar{k}}^\dagger + \alpha_{\bar{k}} \alpha_k \right).
\end{equation}
变分驻点条件要求能量对变分参量的一阶微商为零，等价于消除危险项系数 $H_{20}' \equiv 0$：
\begin{equation}\label{eq:2-3-20}
2\xi_k u_k v_k - \Delta_k (u_k^2 - v_k^2) = 0.
\end{equation}
在正能激发分支 $E_k > 0$ 且 $\Delta_k > 0$ 下，该条件与本任务所选取的准粒子相位约定相容，直接表述为正非对角元的 $2\times 2$ 厄米本征矩阵方程：
\begin{align}\label{eq:2-3-21}
\begin{pmatrix}
\xi_k & +\Delta_k \\
+\Delta_k & -\xi_k
\end{pmatrix}
\begin{pmatrix}
u_k \\
v_k
\end{pmatrix}
= E_k
\begin{pmatrix}
u_k \\
v_k
\end{pmatrix}.
\end{align}
对角化该矩阵给出特征方程 $E_k^2 - (\xi_k^2 + \Delta_k^2) = 0$，解得单准粒子激发能：
\begin{equation}\label{eq:2-3-22}
E_k = \sqrt{\xi_k^2 + \Delta_k^2}.
\end{equation}
第一行本征方程给出 $(\xi_k - E_k) u_k + \Delta_k v_k = 0$。结合逐对归一条件 $u_k^2 + v_k^2 = 1$，直接解出占据与非占据振幅：
\begin{equation}\label{eq:2-3-23}
v_k^2 = \frac{1}{2} \left( 1 - \frac{\xi_k}{E_k} \right), \quad u_k^2 = \frac{1}{2} \left( 1 + \frac{\xi_k}{E_k} \right), \quad u_k v_k = \frac{+\Delta_k}{2 E_k} > 0.
\end{equation}
必须指出：消去 $H_{20}'$ 仅保证能量泛函处于一阶变分驻点（stationarity），该解是否对应局域能量极小，还需检查约束允许方向上的二阶变分及竞争分支。QRPA频率、模范数与零模也要按同一稳定性合同解释，不能只看一组实频率。


\subsubsection{自洽能隙方程与粒子数方程}\label{sec:2-3-自洽能隙方程与粒子数方程}


将乘积项 $u_k v_k = \Delta_k / (2E_k)$ 代回配对能隙与粒子数定义式，给出在全系统中自洽保留自能 $\mu_k = G_{kk} v_k^2$ 的 BCS 方程组：
\begin{equation}\label{eq:2-3-24}
\Delta_k = \sum_{l > 0} G_{kl} \frac{\Delta_l}{2 \sqrt{\xi_l^2 + \Delta_l^2}}, \quad \xi_l = \epsilon_l - \lambda - G_{ll} v_l^2,
\end{equation}
\begin{equation}\label{eq:2-3-25}
\bar{n} = \sum_{k>0} \left( 1 - \frac{\xi_k}{\sqrt{\xi_k^2 + \Delta_k^2}} \right).
\end{equation}
若在活跃壳层内采用恒定配对相互作用力模型（$G_{kl} = G$），则自能为 $\mu_k = G v_k^2$，活跃轨道上的配对能隙统一为常数 $\Delta$，能隙方程严格写为：
\begin{equation}\label{eq:2-3-26}
\frac{2}{G} = \sum_{k>0} \frac{1}{\sqrt{(\epsilon_k - \lambda - G v_k^2)^2 + \Delta^2}}.
\end{equation}
上式约去了共同的非零 $\Delta$，因此不能用它排除或求出 $\Delta=0$ 支；零隙解应回到未约分的自洽方程处理。

在略去高阶准粒子残余相互作用 $V_{\mathrm{RES}}$ 的冻结独立准粒子近似下，两准粒子激发的能量满足 $E_k + E_l \ge 2\Delta$。这一 $2\Delta$ 是共同正能隙、冻结独立准粒子近似下的代数下界，并不要求连续谱；在实际原子核离散单粒子能级结构中，最低双准粒子激发由离散能级和 $E_{k_1} + E_{k_2}$ 决定，并受组态依赖的残余相互作用影响，并不必然等于常数 $2\Delta$。



\subsection{规范对称性破缺、配对转动与物理粒子数守恒}\label{sec:2-3-3}


BCS 理论将多体配对关联映射为非相互作用准粒子均场，但这种简化建立在破缺粒子数规范不变性的基础之上。理解该机制需要明确相关物理量的内涵及真实核体系的物理约束。


\subsubsection{能隙、配对张量与规范角}\label{sec:2-3-能隙配对张量与规范角}


在配对理论中，必须严格区分三个不同范畴的物理概念：
1. \textbf{配对能隙 $\Delta$：} 具有能量量纲（单位通常为 $\text{MeV}$）。它表征破裂一个核子对并激发两个独立准粒子所需的特征能标。
2. \textbf{配对张量振幅 $\kappa_k$：} 无量纲复数振幅，定义为在破缺粒子数参考态下的反常配对密度：
\begin{equation}\label{eq:2-3-27}
\kappa_k \equiv \langle\mathrm{BCS}| c_{\bar{k}} c_k |\mathrm{BCS}\rangle = u_k v_k = \frac{\Delta_k}{2 E_k}.
\end{equation}
   它表征特定轨道对在均场参考态下的凝聚振幅。这里需要区分定粒子数态与破缺参考态：\textbf{在守恒核子数 $N$ 的真实物理本征态 $|N\rangle$ 下，由于规范对称性保持，反常密度期望值恒等于零（$\langle N | c_{\bar{k}} c_k | N \rangle \equiv 0$）。} 这是因为算符 $c_{\bar{k}} c_k$ 改变体系总粒子数 2。非零反常密度用于表征选定规范取向的破缺参考态；在真实原子核态中，超流配对关联由粒子数守恒的两体密度矩阵关联 $\langle N | P_k^\dagger P_l | N \rangle$、双核子转移反应矩阵元 $\langle A+2 | P_k^\dagger | A \rangle$ 以及定数配对关联能来定量表征。
3. \textbf{配对规范角 $\theta$：} 全局粒子数规范变换群 $U(1)$ 的无量纲变换参量。在单粒子规范变换 $c_k \to e^{-i\theta} c_k$ 下，双核子湮灭算符变换为 $c_{\bar{k}} c_k \to e^{-2i\theta} c_{\bar{k}} c_k$。无量纲粒子数算符 $\hat{n}$ 在角表象下对应微分算符 $\hat{n} = -i \partial/\partial\theta$（不含 $\hbar$）。需要强调的是，准粒子无规相位近似（QRPA）中构造的局域共轭角算符是参考态局域切空间上的线性化描述，并不能推广为有限核完整 Hilbert 空间中满足严格全局对易关系的自伴量子力学角算符。


\subsubsection{均场破缺与局域配对转动}\label{sec:2-3-均场破缺与局域配对转动}


孤立原子核体系遵循严格的粒子数守恒定律（$[\hat{H}, \hat{n}] = 0$）。BCS 或 Hartree–Fock–Bogoliubov（HFB）参考态中出现的粒子数规范对称性破缺，是平均场变分近似处理复杂吸引配对相互作用的有效数学手段，不代表物理原子核发生真实的粒子数不守恒。

在自洽均场及残余响应保留相应对称性的理论极限中，破缺的连续规范对称性对应 Nambu–Goldstone（NG）零能模。Hinohara（2015）在线性响应理论框架下探讨了粒子数规范破缺态的零频响应\cite{hinohara-2015-nambu-goldstone-inertia-response}。在总能量 $E = -B_+$ 关于粒子数偏离的二次微商展开中：
\begin{equation}\label{eq:2-3-28}
E(N) = E(N_0) + \lambda \, \delta N + \frac{(\delta N)^2}{2 \mathcal{J}_{\text{pair}}} + \cdots,
\end{equation}
相应的正束缚能 $B_+(N) = -E(N)$ 展开式为：
\begin{equation}\label{eq:2-3-29}
B_+(N) = B_+(N_0) - \lambda \, \delta N - \frac{(\delta N)^2}{2 \mathcal{J}_{\text{pair}}} + \cdots,
\end{equation}
其中化学势为 $\lambda = \partial E/\partial N$，配对转动惯量定义为：
\begin{equation}\label{eq:2-3-30}
\mathcal{J}_{\text{pair}} \equiv \frac{\partial \bar{N}}{\partial \lambda} = -S(\hat{N}, \omega = 0) \quad (\text{单位：}\text{MeV}^{-1}).
\end{equation}
能面二次曲率满足 $\mathcal{J}_{\text{pair}}^{-1} = \partial^2 E/\partial N^2 = -\partial^2 B_+/\partial N^2$。此处始终以E表示内部总能，以$B_+$表示正束缚能；这一局部展开还要求指定平滑的自洽态族、固定其它粒子数的路径及非奇异响应。

必须明确区分两类物理转动惯量：\textbf{配对转动惯量 $\mathcal{J}_{\text{pair}}$ 的量纲是 $\text{MeV}^{-1}$，反映的是能量随粒子数规范变量变化的微观曲率响应；而空间三维集体转动惯量 $\mathcal{J}_{\text{space}}$ 的量纲为 $\hbar^2/\text{MeV}$（或记为 $\text{MeV}^{-1}\hbar^2$），反映的是孤立体系绕真实空间几何轴旋转的动能响应。二者对应的微观自由度与动力学机制完全不同，绝对不能因名称相似而混同。}

为了消除均场近似引入的粒子数弥散，理论上采用粒子数投影（Particle Number Projection, PNP）方法，通过对规范角进行闭合积分：
\begin{equation}\label{eq:2-3-31}
\hat{P}_N = \frac{1}{2\pi} \int_0^{2\pi} e^{i (\hat{n} - N) \theta} \mathrm{d}\theta,
\end{equation}
从试探波函数中滤出具有好粒子数的分量 $|N\rangle \propto \hat{P}_N |\mathrm{BCS}\rangle$。需要指出的是，投影后的状态消除了解中的粒子数涨落，属于好粒子数近似态，但由于其试探波函数形式仍受限于投影前的变分空间，通常并不等同于完整多体哈密顿量的严格本征态。这一内禀参考态与对称性恢复的逻辑，构成了后续处理空间旋转对称性恢复（角动量投影）的方法范式。



\subsection{阻塞效应与多准粒子态的微观重排}\label{sec:2-3-4}


当原子核处于奇 $A$ 核、奇奇核或偶偶核的断对激发态时，体系包含未成对的单粒子。这些未配对核子对周围核子介质的配对散射产生显著的排斥与阻碍作用，即阻塞效应（blocking effect）\cite{book-jain-nuclear-isomers}。


\subsubsection{泡利阻塞与自洽重排}\label{sec:2-3-泡利阻塞与自洽重排}


设某一特定的单核子态 $|k_b\rangle$ 已被一个奇准粒子单独占据。根据费米子泡利不相容原理：
1. 任何其他同类核子均无法再填入该态 $|k_b\rangle$；
2. 原本对应的时间反演轨道对 $(k_b, \bar{k}_b)$ 无法再容纳双核子成对凝聚；
3. 双核子配对散射相互作用中，所有涉及轨道 $k_b$ 的对散射矩阵元被全部阻断。

因此，轨道 $k_b$ 被从活跃的配对相互作用空间中剔除。在自洽计算中，该态的自洽能隙方程与粒子数方程变为：
\begin{equation}\label{eq:2-3-32}
\Delta_k = \sum_{l > 0, l \ne k_b} G_{kl} \frac{\Delta_l}{2 \sqrt{\xi_l^2 + \Delta_l^2}},
\end{equation}
\begin{equation}\label{eq:2-3-33}
\bar{n} - 1 = \sum_{k>0, k \ne k_b} \left( 1 - \frac{\xi_k}{\sqrt{\xi_k^2 + \Delta_k^2}} \right).
\end{equation}
当被阻塞的轨道位于费米面附近时，该能级本应对配对能隙方程提供较大贡献。将其剔除后，自洽求解得到的有效配对能隙 $\Delta$ 出现局域减小。与此同时，奇核子占据引起的势场变化与粒子数守恒要求，会驱动周围非阻塞轨道的化学势 $\lambda$ 与占据几率 $v_k^2$ 发生自洽调整，引起核芯的极化重排效应\cite{book-jain-nuclear-isomers}。


\subsubsection{多准粒子态与高 $K$ 异能态机制}\label{sec:2-3-多准粒子态与高k异能态机制}


在形变核区，当体系具有良好的轴对称性（或接近轴对称）且破坏轴对称性的耦合或 Coriolis 混合较弱时，单粒子角动量在对称轴上的投影量子数 $\Omega$ 近似为守恒量；若核体系显著偏离轴对称（具有显著三轴形变），$\Omega$ 通常发生混合，轴投影不再是良好量子数。在近轴对称条件下断开多个时间反演对时，可以激发多个准粒子，其单粒子轴投影平行极化耦合形成较大的总内禀角动量投影 $K = \sum_i \Omega_i$\cite{book-jain-nuclear-isomers}。

在多准粒子态激发能估算中，通常表示为各准粒子能量、核芯转动能与两体残余作用贡献的叠加\cite{book-jain-nuclear-isomers}：
\begin{equation}\label{eq:2-3-34}
E(K) = E_0 + \sum_{i\in\mathrm{qp}} E_{\mathrm{qp}}(i) + E_{\mathrm{rot}} + \Delta E_{\mathrm{residual}}.
\end{equation}
当多个高投影轨道同时被单粒子占据时，多个配对散射通道被同步阻断，在相对较低的激发能量下形成高 $K$ 构型。需要指出的是，高 $K$ 态形成长寿命同核异能态的关键，不仅在于多准粒子态的能量稳定性，更在于其向低 $K$ 基态带退激时必须跨越巨大的 $K$ 量子数差（$\Delta K$），受到电磁辐射选择定则的极强跃迁阻碍（$K$ 禁戒跃迁阻碍）；阻塞机制解释了多准粒子态在能谱中的形成与局域能隙减弱，而长半衰期的产生还取决于该阻碍机制与退激通道的动力学特征（例如 $^{178}\mathrm{Hf}$ 中的 $16^+$ 态半衰期达 $T_{1/2} \approx 31\text{ 年}$，其具体跃迁矩阵元与寿命分析将在后续章节专门展开）。


\subsubsection{阻塞状态下的配对残留与均场近似边界}\label{sec:2-3-阻塞状态下的配对残留与均场近似边界}


关于多准粒子态与阻塞效应，需要确立客观的理论适用界限：

首先，阻塞若干个特定单粒子轨道虽然削弱了局域配对能隙，但未被阻塞的其他价轨道仍然在费米面两侧参与配对散射。即便在四准粒子或六准粒子激发态中，剩余轨道仍可能维持配对关联，其大小须在指定相互作用和组态下计算。不能断言多准粒子态的微观配对完全消失，也不能将多准粒子态的微观转动惯量直接等同为无配对的理想刚体转动惯量。

其次，静态BCS的非零能隙解可能在临界强度或旋转条件下消失，这不等于粒子数守恒体系的全部配对关联同时消失。下面的两壳例子给出连续趋零的能隙，而其它自洽分支也可能发生交叉或跳变。有限核没有无限体系的热力学极限，仍可有精确能级交叉和选支后的非解析行为。粒子数投影、精确对角化或更广的动力学描述可检验均场遗漏的关联；它们不保证所有路径都平滑，也不能仅由弱配对便断言存在强烈的配对振动。



\subsection{解析控制：两能级教学极限与集体响应的配对反馈}\label{sec:2-3-5}


为了在可控模型下检验均场配对解的临界行为，考察简化的两能级配对模型空间。该模型允许推导精确的解析表达式，用于展示变分解的数学结构及其局限性。


\subsubsection{两能级模型在纯配对哈密顿量下的解析推导}\label{sec:2-3-两能级模型在纯配对哈密顿量下的解析推导}


设模型空间仅包含两个能级，单粒子能量分别为 $\epsilon_1 = -\frac{1}{2}\varepsilon$ 与 $\epsilon_2 = +\frac{1}{2}\varepsilon$，能级间距为 $\varepsilon$。每个壳层的轨道对简并度均为 $\Omega = (2j+1)/2$。设体系的平均粒子数严格等于单个能级的容量，即 $\bar{n} = 2\Omega$（半填充条件）。在无相互作用时，下能级完全被 $2\Omega$ 个核子填满（$v_1 = 1, v_2 = 0$），上能级完全为空。

在本任务采用的纯配对哈密顿量 $\hat{H} = \sum \epsilon_k n_k - G \sum_{k, l} P_k^\dagger P_l$ 下，保留完整的对角项与其 Wick 自能修正 $\mu_i = G v_i^2$。定义上下两能级的占据差参量 $t \equiv v_1^2 - v_2^2$。由粒子—空穴对称性与粒子数守恒要求 $v_1^2 + v_2^2 = 1$，得到 $v_1^2 = (1+t)/2$ 与 $v_2^2 = (1-t)/2$。

在这条镜像对称的BCS积态路径上，直接计算能量期望得到
\begin{equation}\label{eq:2-3-35}
\mathcal E(t)=-\Omega\varepsilon t-G\left[\Omega+\frac{(4\Omega^2-2\Omega)(1-t^2)}4\right].
\end{equation}
其一阶导数为 $-\Omega\varepsilon+G\Omega(2\Omega-1)t$，因此内部驻点满足：
\begin{equation}\label{eq:2-3-36}
t_* = \frac{\varepsilon}{G(2\Omega - 1)}.
\end{equation}
在非零配对且保持镜像对称的驻点上，自洽化学势为 $\lambda = -\frac{G}{2}$，有效单粒子能量满足：
\begin{equation}\label{eq:2-3-37}
|\xi_1| = |\xi_2| = \frac{\Omega \varepsilon}{2\Omega - 1}.
\end{equation}
非零能隙方程在两壳各计 $\Omega$ 对轨道，给出 $1=G\Omega/E$，故准粒子能量为：
\begin{equation}\label{eq:2-3-38}
E_1 = E_2 = E = \Omega G.
\end{equation}
进而解出配对能隙的解析表达式：
\begin{equation}\label{eq:2-3-39}
\Delta = \sqrt{E^2 - \xi^2} = \Omega G \sqrt{1 - t_*^2} = \Omega G \sqrt{1 - \left[ \frac{\varepsilon}{G(2\Omega - 1)} \right]^2} \quad (G > G_{c,\mathrm{pair}}).
\end{equation}
临界配对相互作用强度为：
\begin{equation}\label{eq:2-3-40}
G_{c,\mathrm{pair}} = \frac{\varepsilon}{2\Omega - 1}.
\end{equation}
在非零配对支（$G > G_{c,\mathrm{pair}}$）上，占据几率为：
\begin{equation}\label{eq:2-3-41}
v_1^2 = \frac{1}{2} \left[ 1 + \frac{\varepsilon}{G(2\Omega - 1)} \right], \quad v_2^2 = \frac{1}{2} \left[ 1 - \frac{\varepsilon}{G(2\Omega - 1)} \right].
\end{equation}
此处给出明确说明：Suhonen 专著第 13.5 节（印刷本第 408–410 页）得到的临界强度为 $2\varepsilon/(4\Omega - 1)$，其分母多出的因子源于 Suhonen 所采用的哈密顿量不仅包含时间反演轨道对的单极散射，还包含了所有壳层指标 $a,b,c,d \in \{1,2\}$ 耦合为 $J=0$ 的广义两体矩阵元，从而在自能中引入了额外的跨壳粒子—空穴收缩贡献；完整对账与模型映射留待附录 B。上述解析解是在本节正文明确规定的纯配对相互作用哈密顿量与自能约定下的严格推导结果。


\subsubsection{极限检验与临界条件}\label{sec:2-3-极限检验与临界条件}


利用上述解析表达式可清晰考察变分解的极限性质：
1. \textbf{临界阈值与连续消失：} 当相互作用强度从上方逼近临界点 $G \downarrow G_{c,\mathrm{pair}}$ 时，由于根式行为 $\Delta \sim \Omega G \sqrt{2(G - G_c)/G_c}$，能隙 $\Delta$ 连续趋近于零。当 $G \le G_{c,\mathrm{pair}}$ 时，能隙方程无非平凡实根，$\Delta \equiv 0$，此时占据几率退化为平庸的 $v_1 = 1, v_2 = 0$，体系回到无配对均场的基态。
2. \textbf{区分临界极限与无相互作用极限：} 必须明确区分能隙消失的临界极限与无两体作用极限。当 $G \to G_{c,\mathrm{pair}}^+$ 时，均场配对能隙 $\Delta \to 0$，但此时两体配对相互作用强度 $G_{c,\mathrm{pair}} > 0$ 仍为有限正值，体系在对角部分仍包含两体残余作用；真正的无相互作用极限对应 $G \to 0$。将二者混淆会误判临界区附近的相互作用属性。
3. \textbf{强配对极限（$G/\varepsilon \to \infty$）：} 当 $G/\varepsilon \gg 1$ 时，$t_* \to 0$，占据几率趋于均分 $v_1^2 \to 1/2, v_2^2 \to 1/2$。此时能隙与相互作用强度之比趋于饱和比值 $\Delta / G \to \Omega$。

两能级模型是用于概念检验的高度简化模型，其解析解反映的是特定简并空间与恒定配对矩阵元假定下的数学特征，不能作为实际原子核结构计算的直接结果。 在真实核体系中，单粒子能级呈现不均匀分布与形变劈裂，且残余相互作用包含高阶多极分量，其实际演化受复杂组态混合控制。


\subsubsection{配对对集体转动与动态响应的微观反馈}\label{sec:2-3-配对对集体转动与动态响应的微观反馈}


在微观多体框架下，超流配对关联对核体系的宏观集体运动具有关键的调制作用：
1. \textbf{两准粒子激发阈值的施加：} 准粒子能量满足 $E_k = \sqrt{\xi_k^2 + \Delta_k^2} \ge \Delta_k$，使得费米面附近的内禀两准粒子激发能受到能隙推高。这可以拉开部分集体态与内禀二准粒子组态的能标；实际阻尼还取决于耦合矩阵元与可用相空间；
2. \textbf{跃迁相干因子调制：} 在准粒子表象下展开单体电磁与多极跃迁算符时，由于配对产生的粒子与空穴成分之间的量子干涉，跃迁矩阵元乘上了相干因子 $(u_k v_{k'} \mp v_k u_{k'})$\cite{book-suhonen-from-nucleons-to-nucleus}。这导致不同多极跃迁的响应截然不同，部分多极跃迁受到相干增强，而另一些则受到抑制；
3. \textbf{转动惯量的微观抑制：} 在 Inglis–Belyaev 显式曲率公式中，转动惯量分量由角动量算符在准粒子基底下的矩阵元除以双准粒子激发能给出，其表达式受两准粒子能量分母 $E_k + E_{k'}$ 增大与相干因子干涉的双重压低，能够改变转动惯量；无配对极限仍保留壳结构与响应效应，并不自动等于经典刚体值。

然而，在动力学层面必须明确：\textbf{集体转动惯量是单粒子轨道能级、配对能隙、形状形变度以及多极残余相互作用联合贡献的微观响应结果，改变配对能隙并不单一决定转动惯量的变化；在没有完整自洽微观动力学计算的前提下，不能简单将转动惯量的偏离唯一归因于配对的消失。} 这些概念为后续章节讨论三维旋转、微观惯量张量与形状动力学提供了自洽的微观出发点。
